\begin{entry}{Q020}{What survives a Haar-random gate}

\entrymeta{
  topic   = {Haar measure; random matrices; unitary twirl.},
  est     = {$\sim$15 minutes},
  origin  = {Original.},
  source  = {Original problem.},
  prereq  = {None beyond linear algebra; the two averaging rules are stated below.},
  status  = drafted,
}

\begin{problem}
\textbf{What this describes.} A Haar-random unitary is the idealisation of a maximally
generic, chaotic evolution: no direction in Hilbert space is preferred. Averaging over it
is how random circuits are analysed, and the basic question is what information about an
input survives the average. Everything follows from two rules.

\textbf{Setup.} A random matrix is a matrix whose entries are random numbers. Its average
$\mathbb E[M]$ is the matrix of entry-wise averages, $(\mathbb E[M])_{kl}=\mathbb E[M_{kl}]$.

\emph{Rule 1 (linearity).} For fixed matrices $A,B$ and a random matrix $M$,
\[
  \mathbb E[A\,M\,B]=A\;\mathbb E[M]\;B .
\]
Indeed $(AMB)_{kl}=\sum_{i,j}A_{ki}M_{ij}B_{jl}$; the average of a sum is the sum of the
averages, and constants factor out.

\emph{Rule 2 (invariance).} If $U$ is Haar-distributed on $\mathrm U(n)$ and $V$ is a
fixed unitary, then $VU$ has the same distribution as $U$, so for any matrix-valued
function $f$,
\[
  \mathbb E[f(VU)]=\mathbb E[f(U)] .
\]

Let $U$ be Haar-distributed on $\mathrm U(n)$, $X$ a fixed $n\times n$ matrix, and
$M:=\mathbb E[UXU^{\dagger}]$.

\begin{enumerate}[label=\textbf{(\alph*)},leftmargin=*]
  \item \emph{Warm-up.} Use the two rules with a well-chosen $V$ to show that
        $\mathbb E[U]=0$.

  \item \emph{Invariance of $M$.} Apply Rule~2 to $f(U)=UXU^{\dagger}$, then Rule~1,
        to show that
        \begin{equation}
            VMV^{\dagger}=M\qquad\text{for every unitary } V 
            \label{invariance identity}
        \end{equation}
          

  \item \emph{$M$ is a multiple of the identity.} Use (b) with two families of $V$:
        \begin{itemize}
          \item $V$ diagonal, with one entry equal to $-1$ and all others $+1$. What does
                this force on the off-diagonal entries of $M$?
          \item $V$ the permutation swapping two basis vectors and fixing the rest. What
                does this force on the diagonal entries?
        \end{itemize}
        Conclude that $M=c\,\id$ for some $c\in\CC$.

  \item \emph{The constant.} Take the trace of $M=\mathbb E[UXU^{\dagger}]$; since the
        trace is a sum of entries, it can be exchanged with the average. Find $c$.

  \item \emph{Reading.} Let $\rho$ be any state of two qubits, so $n=4$. What is the
        average state after a Haar-random gate, and what property of $\rho$ survives?
\end{enumerate}
\end{problem}

% ============================================================== solutions ===

\begin{solution}[Part (a)]
For any fixed unitary $V$, by linearity
\begin{equation*}
    \mathbb E[VU] = V\mathbb E[U].
\end{equation*}
Now, take $f(X)=X: \CC^{n\times n} \rightarrow \CC^{n\times n}$, so by invariance
\begin{equation*}
    \mathbb E[VU] = \mathbb E[U].
\end{equation*}
Thus, taking $V=-\id$, $ \mathbb E[U]=0$.

\end{solution}

\begin{solution}[Part (b)]
Given a unitary $V$, by the invariance $\mathbb E(VUXU^\dagger V^\dagger)=\mathbb E(UXU^\dagger)=M$, using linearity on the left hand side $M = V \mathbb E(UXU^\dagger)V^\dagger=V  MV^\dagger$.
\end{solution}

\begin{solution}[Part (c)]
Let $V_{ij}= (-1)^{\delta_{ia}}\delta_{ij}$, so the $a-$th element of the diagonal is the $-1$. Then
\begin{equation}
   M_{kl}= \sum_{i,j} (-1)^{\delta_{ka}}\delta_{ki}M_{ij} (-1)^{\delta_{ja}}\delta_{jl} = (-1)^{\delta_{ka}+\delta_{ja}} M_{lk}
\end{equation}
So for diagonal elements, this equation tells nothing. However, given a pair $(l,k)$ of off-diagonal elements $l\neq k$, one can take $a=k$ (due to the generality of $V$), leading to $M_{lk}=-M_{kl}$. Thus, the matrix $M$ must be diagonal.

In order to the unitary act as a permutation of two basis vectors, it must act on that subspace as the Pauli-X matrix, $ \ket{e_i}\bra{e_j} + \ket{e_j}\bra{e_i}$, and act trivially as the identity on the rest 
$$V=\ket{e_i}\bra{e_j} + \ket{e_j}\bra{e_i}+ \sum_{k\neq i,j} \ket{e_k}\bra{e_k}.$$ 
This forces $(VMV^\dagger)_{ii}=M_{jj}$ and $(VMV^\dagger)_{jj}=M_{ii}$, while leaving the others diagonal elements unaltered. 

Furthermore, Eq.\eqref{invariance identity} imposes $M_{jj}=M_{ii}$ for all diagonal elements, such that $M$ must be proportional to the identity.

\end{solution}

\begin{solution}[Part (d)]
Using linearity and cyclicity of the trace 
\begin{equation}
    c \id= E[UXU^{\dagger}]   \implies c\, n= \Tr\left( \mathbb E[UXU^{\dagger}]\right) = E[\Tr(UXU^{\dagger})] = \mathbb E[ \Tr(X)]= \Tr(X)
\end{equation}
\end{solution}

\begin{solution}[Part (e)]
For any state $\mathbb E(U\rho U^\dagger)= \frac{1}{n} \id$, which is the maximally mixed state. Meaning the operation erases all information from the state. 
\end{solution}



\attempt{
  date      = 2026-10-05,
  time      = 20min,
  difficulty= 1,
  howitwent = ,
  insight   = ,
}

\end{entry}